\chapter{协程入门：异步流程不是线程}

\section{问题入口：回调嵌套为什么难维护}

异步 I/O 常见写法是回调：

\begin{lstlisting}[language=C++,caption={回调嵌套示意}]
read_request(socket, [socket](Request request) {
    query_database(request.user_id, [socket](User user) {
        write_response(socket, make_response(user));
    });
});
\end{lstlisting}

逻辑顺序是“读请求、查数据库、写响应”，但代码结构变成嵌套。错误处理、取消和资源释放会更复杂。协程提供一种把异步流程写得像顺序代码的机制。

\section{协程不是自动多线程}

\cpp20 协程是语言机制，不等于线程池，也不等于事件循环。协程可以暂停和恢复，但谁来恢复它，取决于 awaiter、调度器或库。没有运行时库，语言本身不会替你完成异步 I/O。

伪代码形态：

\begin{lstlisting}[language=C++,caption={协程式异步流程示意}]
Task<Response> handle(Socket socket)
{
    Request request = co_await read_request(socket);
    User user = co_await query_database(request.user_id);
    co_return make_response(user);
}
\end{lstlisting}

这段代码表达的是控制流形态，不是本书要实现完整协程库。学习重点是理解 \code{co_await}、\code{co_return} 和生命周期边界。

\section{co\_await 的直觉}

\code{co_await something} 大致表示：如果结果还没准备好，当前协程可以暂停，把控制权还给调度器；等结果准备好，再从这里继续。它不是阻塞当前线程等待。阻塞会占住线程，暂停协程则可以让线程去做别的工作。

这也是协程适合 I/O 密集任务的原因：大量请求都在等待网络或磁盘时，不需要为每个请求占一个线程。

\section{协程帧和生命周期}

协程暂停时，局部变量必须保存起来。编译器会把协程的状态放进协程帧。于是协程里的引用、指针和视图仍然要遵守生命周期规则：

\begin{lstlisting}[language=C++,caption={协程中捕获引用的风险示意}]
Task<void> bad(std::string_view text)
{
    co_await delay();
    use(text); // text 指向的字符串可能已经销毁
}
\end{lstlisting}

如果协程会跨越调用者作用域继续运行，就不要保存指向调用者临时数据的视图。需要长期使用的数据应该复制或转移所有权。

\section{取消和异常}

协程流程看起来像同步代码，但错误边界仍然要设计。异步操作可能失败，\code{co_await} 处可能抛异常或返回错误。取消也不是强制销毁一切，而是协作协议：正在等待的操作要能收到取消请求，并让协程走到清理路径。

伪代码：

\begin{lstlisting}[language=C++,caption={协程中的错误边界示意}]
Task<void> serve_one(Socket socket)
{
    try {
        Request request = co_await read_request(socket);
        Response response = co_await build_response(request);
        co_await write_response(socket, response);
    } catch (const std::exception& error) {
        log_error(error.what());
    }
}
\end{lstlisting}

捕获异常必须有目的。这里是在请求边界记录错误，防止一个请求失败影响整个服务循环。

\section{什么时候不需要协程}

协程不是所有异步问题的默认答案。下面场景普通代码更合适：

\begin{itemize}
  \item 简单后台任务，用 \code{std::jthread} 足够。
  \item CPU 密集并行，用线程池和任务分解更直接。
  \item 项目没有成熟协程库或调度器。
  \item 团队还无法调试协程栈和生命周期问题。
\end{itemize}

协程适合大量异步等待、流程嵌套严重、已有可靠异步运行时的场景。

\section{从回调到状态机}

即使不用协程，也可以把异步流程看成状态机：

\begin{lstlisting}[caption={请求处理状态}]
ReadRequest -> QueryUser -> WriteResponse -> Done
                          -> Failed
\end{lstlisting}

协程的一个价值，是让编译器帮你保存状态机的局部状态。理解这一点，就不会把协程误解成神秘语法。它是在帮你表达“可暂停的函数”。

\begin{exercisebox}
把“读文件、解析配置、启动服务”写成三步伪异步流程。先用回调写，再用协程伪代码写。比较错误处理放在哪里更清楚。重点不是编译运行，而是理解控制流形状。
\end{exercisebox}

\section{协程背后的手写状态机}

协程最容易被误解成魔法。先把这段伪代码拆开：

\begin{lstlisting}[language=C++,caption={协程顺序写法}]
Task<Response> handle(Socket socket)
{
    Request request = co_await read_request(socket);
    User user = co_await query_database(request.user_id);
    co_return make_response(user);
}
\end{lstlisting}

如果不用协程，可以手写一个状态机：

\begin{lstlisting}[caption={手写状态机形态}]
state = ReadRequest;
保存 socket;

当 read_request 完成:
    保存 request;
    state = QueryDatabase;
    发起 query_database;

当 query_database 完成:
    保存 user;
    state = Done;
    生成 response;
\end{lstlisting}

协程的作用，是让编译器帮你生成和维护这套状态机。局部变量 \code{request}、\code{user}、\code{socket} 如果跨越暂停点还要使用，就会被保存到协程帧里。所谓协程帧，可以理解成“堆上或运行时管理的一块状态机存储”。

\section{co await 近似展开}

\code{co_await read_request(socket)} 大致做三件事：

\begin{enumerate}
  \item 问 awaiter：结果是否已经准备好。
  \item 如果没准备好，保存当前协程状态并把控制权还给调度器。
  \item 等外部事件完成后，调度器恢复协程，从暂停点继续。
\end{enumerate}

这和阻塞调用不同。阻塞调用会让当前线程卡住；协程暂停通常让线程回到事件循环或调度器，去处理别的任务。具体行为由库实现决定，语言只提供暂停和恢复的协议。

\section{生命周期问题为什么更隐蔽}

普通函数返回后，局部变量销毁，调用栈消失。协程暂停后，局部状态可能还活在协程帧里。于是下面的函数看起来只是接收一个 \code{string_view}，但暂停后风险变大：

\begin{lstlisting}[language=C++,caption={跨暂停点保存视图}]
Task<void> log_later(std::string_view message)
{
    co_await delay();
    write_log(message);
}
\end{lstlisting}

如果调用者传入临时字符串，暂停期间临时对象早就销毁。更稳的接口是拥有数据：

\begin{lstlisting}[language=C++,caption={协程跨暂停点拥有字符串}]
Task<void> log_later(std::string message)
{
    co_await delay();
    write_log(message);
}
\end{lstlisting}

这不是说协程里不能用视图，而是要看视图是否跨越暂停点。跨暂停点的数据，生命周期要求要按异步任务处理。

\section{协程库边界}

C++20 只给语言机制，不给完整任务类型。教材里的 \code{Task<T>} 是伪类型，真实项目需要库来定义：

\begin{itemize}
  \item 协程如何启动。
  \item 暂停后由谁保存句柄。
  \item I/O 完成后由谁恢复协程。
  \item 异常如何存储和传播。
  \item 取消如何通知等待中的操作。
\end{itemize}

所以学习协程时要分清两层：语言层的 \code{co_await}、\code{co_return}、协程帧；库层的 \code{Task}、调度器、事件循环和 I/O awaiter。没有库层，语言层不能单独完成服务器。

\section{把回调错误处理改成协程错误边界}

回调版本里，错误处理容易散落：

\begin{lstlisting}[language=C++,caption={回调错误散落}]
read_request(socket,
    [socket](Request request) {
        query_database(request.user_id,
            [socket](User user) {
                write_response(socket, make_response(user));
            },
            [](Error error) {
                log_error(error);
            });
    },
    [](Error error) {
        log_error(error);
    });
\end{lstlisting}

协程版本可以把一个请求的错误边界集中起来：

\begin{lstlisting}[language=C++,caption={协程请求边界}]
Task<void> serve_one(Socket socket)
{
    try {
        Request request = co_await read_request(socket);
        User user = co_await query_database(request.user_id);
        co_await write_response(socket, make_response(user));
    } catch (const RequestError& error) {
        log_request_error(error);
    }
}
\end{lstlisting}

这就是协程的实际价值：不是自动更快，而是把异步流程恢复成接近顺序代码的形状，让错误、取消和资源释放更容易放在同一个边界里。
